% ======================================================================
%  From the sharp-interface homogenization model of Calonne et al. to an
%  equivalent phase-field formulation for effective thermal conductivity.
%
%  Notation follows Moure & Fu (2024) where the two overlap; Calonne
%  et al. (2011) symbols are mapped onto it. Self-contained.
%  Requires only amsmath, amssymb, amsthm, booktabs, geometry.
% ======================================================================

\documentclass[11pt]{article}
\usepackage{amsmath,amssymb,amsthm,booktabs}
\usepackage[margin=1in]{geometry}

\newtheorem{lemma}{Lemma}
\newtheorem{proposition}{Proposition}
\newtheorem{corollary}{Corollary}
\newtheorem{remark}{Remark}

% --- notation macros (Moure & Fu-aligned) ---
\newcommand{\Ki}{K_i}
\newcommand{\Ka}{K_a}
\newcommand{\Keff}{\mathbf{K}_{\mathrm{eff}}}
\newcommand{\phibar}{\overline{\phi}}
\newcommand{\phii}{\phi_i}
\newcommand{\phia}{\phi_a}
\newcommand{\Om}{\Omega}
\newcommand{\Omi}{\Omega_i}
\newcommand{\Oma}{\Omega_a}
\newcommand{\Gam}{\Gamma}
\newcommand{\bn}{\hat{\mathbf n}}
\newcommand{\bt}{\mathbf t}
\newcommand{\bI}{\mathbf I}
\newcommand{\be}{\mathbf e}
\newcommand{\jump}[1]{[\![\,#1\,]\!]}
\newcommand{\Hper}{H^1_{\mathrm{per}}(\Om)}
\newcommand{\ind}[1]{\mathbf 1_{#1}}
\newcommand{\chiI}{\chi}
\newcommand{\dK}{\delta\! K}
\newcommand{\Band}{B_\varepsilon}
\newcommand{\avg}[1]{\big\langle #1 \big\rangle}

\begin{document}

\begin{center}
{\large\bfseries From Calonne et al.'s sharp-interface homogenization\\[2pt]
to an equivalent phase-field formulation of $\Keff$}
\end{center}

\section{What we are computing and why}

Snow, firn, and icy regolith are two-phase composites of ice and pore air.
Heat conducts through \emph{both} phases: the ice conductivity $\Ki$ and the
air conductivity $\Ka$ differ by nearly two orders of magnitude
($\Ki/\Ka \approx 2.29/0.02 \approx 114$), so neither can be ignored. The
quantity we want is the \emph{effective} (bulk) thermal conductivity tensor
$\Keff$ that a macroscopic sample presents, given only the pore-scale geometry.
It is defined so that the macroscopic heat flux obeys Fourier's law
$\mathbf F = -\Keff\,\nabla T$.

Because the pore scale $\ell$ (grains, necks: tens to hundreds of microns) is
far smaller than any macroscopic scale $L$ of interest, the two scales
separate and $\Keff$ can be obtained by \emph{homogenization}: solve one
auxiliary boundary-value problem, the \emph{cell problem}, on a single
representative elementary volume (REV) $\Om$, and average. Calonne et al.\
(2011) do this on a \emph{sharp} (voxelized) ice--air geometry. Our
phase-field model represents the same geometry with a smooth order parameter.

The purpose of this note is to show that the phase-field cell problem we solve
is the \emph{same} boundary-value problem as Calonne's, written on one domain
instead of two. We do this in two steps, and it is worth saying up front that
\emph{neither step is a limit}:

\begin{itemize}
\item[\S4.1] The two-domain sharp problem and the single global weak form are
the same problem. The interface conditions are not discarded; they are absorbed,
one into the solution space and one into the weak form. The step is reversible.
\item[\S4.2] The sharp coefficient $K_\star$ and the phase-field interpolation
$K(\phibar)$ are the \emph{same function}. They are not two models that agree in
a limit; they are one constitutive law evaluated at two different arguments.
\end{itemize}

\noindent
What is left over after those two steps is therefore not a modelling
discrepancy at all, but a purely \emph{geometric} one: the phase field
$\phii^\varepsilon$ is a smoothed version of the sharp indicator. Sections
\ref{sec:diff}--\ref{sec:laminate} quantify exactly that difference, with a
computable error bound at finite $\varepsilon$ and a closed-form answer for a
laminate. No appeal to an unproved convergence theorem is needed anywhere.

\paragraph{What the cell problem means physically.}
Impose a unit macroscopic temperature gradient in a coordinate direction
$\be_j$. The microstructure bends the local temperature field away from a
straight ramp; the \emph{corrector} $t^j$ is exactly that local deviation.
Once we know how the field bends, we average the resulting local heat flux
over the REV to read off the $j$-th column of $\Keff$. There are $d$ corrector
fields (one per direction), collected into the vector $\bt=(t^1,\dots,t^d)$.

\section{Notation}

\begin{center}
\begin{tabular}{@{}lll@{}}
\toprule
Symbol & Meaning & Source \\
\midrule
$\phii,\ \phia$ & ice / air phase-field variables, $\phia=1-\phii$ & Moure \& Fu \\
$\phibar=\{\phii,\phia\}$ & the phase-field state (set of variables) & Moure \& Fu \\
$\varepsilon$ & phase-field interface decay length & Moure \& Fu \\
$\chiI=\ind{\Omi}$ & \emph{indicator function} of the ice subdomain & --- \\
$\Ki,\ \Ka$ & ice / air thermal conductivities & Calonne \\
$K(\phibar)=\Ka+(\Ki-\Ka)\phii$ & interpolated conductivity (arith., Eq.~13 M\&F) & Moure \& Fu \\
$K_\star=\Ki\ind{\Omi}+\Ka\ind{\Oma}$ & sharp (piecewise-constant) coefficient & Calonne \\
$\dK=K(\phii^\varepsilon)-K_\star$ & coefficient discrepancy & --- \\
$\Band=\{0<\phii^\varepsilon<1\}$ & the diffuse band & --- \\
$T$ & temperature & both \\
$\Om,\ \Omi,\ \Oma$ & REV cell, ice subdomain, air subdomain & (Calonne: $W$) \\
$\Gam,\ \bn$ & ice--air interface, unit normal (ice$\to$air) & Calonne \\
$\bt=(t^1,\dots,t^d)$ & corrector (vector); $\nabla\bt+\bI$ is a tensor & Calonne ($\mathbf t$)\\
$\Keff$ & effective conductivity tensor & both \\
$\avg{\cdot}=\frac{1}{|\Om|}\int_\Om(\cdot)$ & REV average & --- \\
\bottomrule
\end{tabular}
\end{center}

\paragraph{The indicator function.}
$\ind{A}$ denotes the \emph{indicator} (or \emph{characteristic}) function of a
set $A$: the function that is $1$ inside $A$ and $0$ outside,
\begin{equation}
\ind{\Omi}(\mathbf x)=
\begin{cases}
1, & \mathbf x\in\Omi \quad(\text{ice}),\\
0, & \mathbf x\in\Oma \quad(\text{pore air}).
\end{cases}
\label{eq:indicator}
\end{equation}
It is the sharp, discontinuous ancestor of the phase field: $\phii$ is a
smoothed $\ind{\Omi}$, taking the same values $1$ and $0$ in the bulk phases but
sweeping continuously between them across a band of decay length $\varepsilon$.
Writing $K_\star=\Ki\ind{\Omi}+\Ka\ind{\Oma}$ is just a compact way of saying
``$K_\star$ equals $\Ki$ in the ice and $\Ka$ in the air.'' We write
$\chiI\equiv\ind{\Omi}$ below where it shortens an equation.

\paragraph{Two arguments, one law.} Throughout, $\phii$ and $\ind{\Omi}$ are two
\emph{arguments} that may be fed to the single conductivity law $K(\cdot)$; they
are not two different laws. Lemma~\ref{lem:identity} is precisely this statement.

\paragraph{One deliberate choice of symbol.}
We never use a single symbol for both the interface width and the
homogenization scale ratio $\eta=\ell/L$. Only $\varepsilon$ appears below;
$\eta\to0$ is what licenses homogenization in the first place and plays no
further role here.

\section{The sharp-interface model (Calonne et al.)}

At the pore scale, steady conduction with no sources reads
\begin{align}
\nabla\!\cdot\!(\Ki\nabla T)&=0 \ \text{in }\Omi, &
\nabla\!\cdot\!(\Ka\nabla T)&=0 \ \text{in }\Oma,
\label{eq:bulk}\\[2pt]
\jump{T}&=0 \ \text{on }\Gam, &
\jump{K\nabla T\!\cdot\!\bn}&=0 \ \text{on }\Gam.
\label{eq:iface}
\end{align}
The two interface conditions \eqref{eq:iface} are the physics: temperature is
\emph{continuous} across the ice--air boundary, and the normal heat flux is
\emph{conserved} (no energy is created or stored at the interface -- perfect
thermal contact, no interfacial resistance).

Homogenization of \eqref{eq:bulk}--\eqref{eq:iface} yields the macroscopic law
$\nabla\!\cdot\!(\Keff\nabla T)=0$, with $\Keff$ obtained from the
\emph{corrector (cell) problem}: find the cell-periodic $\bt$ with
\begin{align}
\nabla\!\cdot\!\big(\Ki(\nabla\bt+\bI)\big)&=\mathbf 0 \ \text{in }\Omi, &
\nabla\!\cdot\!\big(\Ka(\nabla\bt+\bI)\big)&=\mathbf 0 \ \text{in }\Oma,
\label{eq:cellbulk}\\[2pt]
\jump{\bt}&=\mathbf 0 \ \text{on }\Gam, &
\jump{K(\nabla\bt+\bI)\!\cdot\!\bn}&=\mathbf 0 \ \text{on }\Gam,
\label{eq:celliface}
\end{align}
subject to the zero-mean normalization $\avg{\bt}=\mathbf 0$
(which removes the arbitrary additive constant), and then
\begin{equation}
\Keff=\frac{1}{|\Om|}\left[\int_{\Omi}\!\Ki(\nabla\bt+\bI)\,d\Om
                          +\int_{\Oma}\!\Ka(\nabla\bt+\bI)\,d\Om\right].
\label{eq:keffsharp}
\end{equation}
Equations \eqref{eq:cellbulk}--\eqref{eq:keffsharp} are Calonne et al.\ (2011),
Eqs.~(7)--(11). The corrector problem carries the \emph{same} interface
conditions \eqref{eq:celliface} as the physical problem, for the same reason:
$\bt$ is a temperature-like field, so it is continuous and its flux is
conserved.

\section{Construction of the phase-field equivalent}

We now build the phase-field cell problem \emph{from}
\eqref{eq:cellbulk}--\eqref{eq:keffsharp} in two steps. The geometry is frozen:
$\phii(\mathbf x)$ is a fixed field, $=1$ in ice, $=0$ in pore, transitioning
smoothly over decay length $\varepsilon$; the sharp indicator
\eqref{eq:indicator} it regularizes is $\ind{\Omi}$.

\subsection{Step 1 --- collapse the two domains into one weak statement}

Multiply each bulk equation in \eqref{eq:cellbulk} by an arbitrary
cell-periodic test function $v$, integrate over the respective subdomain, and
integrate by parts. The outer boundary terms cancel by periodicity, leaving an
interface term:
\begin{equation}
0=-\int_{\Om}\nabla v\cdot K_\star(\nabla\bt+\bI)\,d\Om
 \;+\;\int_{\Gam} v\,\jump{K(\nabla\bt+\bI)\!\cdot\!\bn}\,d\Gam.
\label{eq:glue}
\end{equation}
The interface integral vanishes \emph{identically} by the flux
condition in \eqref{eq:celliface}. What remains is a single equation over the
whole cell,
\begin{equation}
\int_{\Om}\nabla v\cdot K_\star(\nabla\bt+\bI)\,d\Om=0
\qquad\forall\,v\in\Hper,
\label{eq:weaksharp}
\end{equation}
and it is well posed only because $\bt$ is a \emph{single} $H^1$ field with no
jump --- which is precisely the continuity condition in
\eqref{eq:celliface}. Thus the two interface conditions have not been
discarded: continuity $\jump{\bt}=0$ is absorbed into the choice of solution
space (a single continuous field), and flux conservation
$\jump{K(\nabla\bt+\bI)\!\cdot\!\bn}=0$ is absorbed into the weak form (it is
the \emph{natural} condition, satisfied automatically by every weak solution).
The step is reversible: restricting $v$ to each subdomain recovers
\eqref{eq:cellbulk}, and undoing the integration by parts recovers
\eqref{eq:celliface}. So \eqref{eq:weaksharp} is equivalent to
\eqref{eq:cellbulk}--\eqref{eq:celliface} --- the sharp two-domain problem and
the single global weak form are the same problem.

\subsection{Step 2 --- the coefficient is already the phase-field interpolation}
\label{sec:step2}

It is tempting to describe the next step as ``replace the sharp coefficient
$K_\star$ by a smooth interpolation that converges to it.'' That description is
weaker than the truth and invites the objection that an unproved limit is being
assumed. In fact no replacement occurs, because the two coefficients are
the same function of the phase variable.

\begin{lemma}[The sharp coefficient \emph{is} the phase-field interpolation]
\label{lem:identity}
Let $K(\phibar)$ be any conductivity law satisfying the two endpoint conditions
\begin{equation}
K(\phii{=}0)=\Ka,\qquad K(\phii{=}1)=\Ki .
\label{eq:endpoints}
\end{equation}
Then, evaluated on the sharp indicator,
\begin{equation}
K(\phibar)\Big|_{\phii=\ind{\Omi}} \;=\; K_\star
\qquad\text{pointwise in }\Om,
\label{eq:identity}
\end{equation}
exactly. No limit, and no smallness of $\varepsilon$, is involved.
\end{lemma}

\begin{proof}
$\Om=\Omi\cup\Oma$ is a disjoint decomposition (up to the measure-zero surface
$\Gam$), so $\ind{\Oma}=1-\ind{\Omi}$ and the indicator $\ind{\Omi}$ takes
\emph{only} the two values $0$ and $1$. At a point of $\Omi$, $\phii=1$, and
\eqref{eq:endpoints} gives $K=\Ki=K_\star$; at a point of $\Oma$, $\phii=0$, and
\eqref{eq:endpoints} gives $K=\Ka=K_\star$. Those are all the points.
\end{proof}

\noindent
For the arithmetic law actually used (Moure \& Fu, Eq.~13) the same statement is
three lines of algebra, which is worth writing out because it is the whole of
Step 2:
\begin{equation}
K_\star
=\Ki\ind{\Omi}+\Ka\ind{\Oma}
=\Ki\ind{\Omi}+\Ka\big(1-\ind{\Omi}\big)
=\Ka+(\Ki-\Ka)\,\ind{\Omi}
=K(\phibar)\Big|_{\phii=\ind{\Omi}} .
\label{eq:algebra}
\end{equation}

\paragraph{What Lemma~\ref{lem:identity} means.}
The sharp model and the phase-field model do not have different constitutive
laws that happen to agree asymptotically. They have \emph{one} law,
$K(\cdot)$, evaluated at two different arguments:
\[
\text{Calonne: } K\big(\ind{\Omi}\big),
\qquad\qquad
\text{phase field: } K\big(\phii^{\varepsilon}\big).
\]
So Step 2 is not an approximation of the \emph{equations}; it is a
regularization of the \emph{geometry}. Every difference between the two models
traces back to a single object, $\phii^\varepsilon-\ind{\Omi}$, which is
supported on the diffuse band and is a statement about how the ice--air boundary
is represented, not about how heat conducts.

A second consequence of Lemma~\ref{lem:identity} is often missed: because the
indicator never takes an intermediate value, the \emph{shape} of the
interpolation between $0$ and $1$ is completely invisible to the sharp problem.
Arithmetic, harmonic, and every other admissible interpolation coincide exactly
on $\ind{\Omi}$. This is why the choice of interpolation can matter only through
the band --- a point we make quantitative in \S\ref{sec:excess}.

\paragraph{The resulting phase-field cell problem.}
Substituting $K(\phibar)$ into \eqref{eq:weaksharp} and \eqref{eq:keffsharp},
and returning the weak form to strong form, gives
\begin{equation}
\boxed{\;
\nabla\!\cdot\!\big(K(\phibar)(\nabla\bt+\bI)\big)=\mathbf 0 \ \text{in }\Om,
\qquad \bt\ \text{cell-periodic},\qquad \int_\Om\bt\,d\Om=\mathbf 0,\;}
\label{eq:cellPF}
\end{equation}
\begin{equation}
\boxed{\;
\Keff=\frac{1}{|\Om|}\int_{\Om} K(\phibar)(\nabla\bt+\bI)\,d\Om.\;}
\label{eq:keffPF}
\end{equation}
Equations \eqref{eq:cellPF}--\eqref{eq:keffPF} are exactly the phase-field
homogenization equations we solve. The two subdomain integrals of
\eqref{eq:keffsharp} have merged into one integral over $\Om$, and the two
interface conditions no longer appear because Step~1 showed they are enforced
implicitly by the single-field weak form. The companion macroscopic equation
$\nabla\!\cdot\!(K(\phibar)\nabla T)=0$ is the same collapse applied to
\eqref{eq:bulk}--\eqref{eq:iface}.

\begin{proposition}[Exact equivalence on a sharp geometry]
\label{prop:exact}
If the phase field is the sharp indicator, $\phii=\ind{\Omi}$, then the
phase-field problem \eqref{eq:cellPF}--\eqref{eq:keffPF} and Calonne's problem
\eqref{eq:cellbulk}--\eqref{eq:keffsharp} are the same boundary-value problem
and return the same $\Keff$ --- identically, not in a limit.
\end{proposition}

\begin{proof}
Step 1 shows \eqref{eq:cellbulk}--\eqref{eq:celliface} is equivalent to
\eqref{eq:weaksharp}; Lemma~\ref{lem:identity} shows the coefficient in
\eqref{eq:weaksharp} is $K(\phibar)|_{\phii=\ind{\Omi}}$, which is the
coefficient in \eqref{eq:cellPF}. The two functionals \eqref{eq:keffsharp} and
\eqref{eq:keffPF} likewise coincide term by term because
$\int_{\Omi}\Ki(\cdot)+\int_{\Oma}\Ka(\cdot)=\int_{\Om}K_\star(\cdot)$.
\end{proof}

\section{Exactly what the two problems differ by}
\label{sec:diff}

At the finite $\varepsilon$ we actually run, $\phii^\varepsilon\neq\ind{\Omi}$,
and Proposition~\ref{prop:exact} no longer applies verbatim. This section says
precisely --- and computably --- what the residual difference is. Nothing here
is asymptotic.

\subsection{The coefficient discrepancy}

For the arithmetic law, $K$ is \emph{affine} in $\phii$, so subtracting
\eqref{eq:algebra} from $K(\phii^\varepsilon)$ gives the discrepancy in closed
form:
\begin{equation}
\dK \;:=\; K(\phii^{\varepsilon})-K_\star
\;=\;(\Ki-\Ka)\,\big(\phii^{\varepsilon}-\chiI\big),
\qquad \operatorname{supp}\dK=\Band .
\label{eq:dK}
\end{equation}
The coefficient error is \emph{exactly proportional} to the geometric error,
with the constant $\Ki-\Ka$, and it lives only where the phase field is strictly
between $0$ and $1$. Both problems are uniformly elliptic on the same bounds,
$\min(\Ka,\Ki)\le K\le\max(\Ka,\Ki)$, since $\phii^\varepsilon\in[0,1]$; so both
are well posed and admit the variational characterization used next.

\subsection{A two-sided bound, with no convergence theorem}

For a symmetric coefficient the cell problem has a Dirichlet minimum principle:
for any unit direction $\be$,
\begin{equation}
A(K)\;:=\;\be\cdot\Keff[K]\,\be
\;=\;\min_{v\in\Hper}\;\frac{1}{|\Om|}\int_\Om(\nabla v+\be)\cdot K\,(\nabla v+\be)\,d\Om
\;=:\;\min_{v\in\Hper} E_K[v].
\label{eq:dirichlet}
\end{equation}
Write $K_1=K(\phii^\varepsilon)$ and $K_2=K_\star$, with minimizers $v_1,v_2$,
and note from \eqref{eq:dK} that
$E_{K_1}[v]-E_{K_2}[v]=\avg{\dK\,|\nabla v+\be|^2}$ for \emph{any} $v$. Testing
each minimizer against the other's energy,
\begin{align}
A(K_1)=E_{K_1}[v_1]&\le E_{K_1}[v_2]=A(K_2)+\avg{\dK\,|\nabla v_2+\be|^2},\\
A(K_2)=E_{K_2}[v_2]&\le E_{K_2}[v_1]=A(K_1)-\avg{\dK\,|\nabla v_1+\be|^2},
\end{align}
which rearrange into a rigorous sandwich on the error:
\begin{equation}
\boxed{\;
\avg{\dK\,|\nabla v_\varepsilon+\be|^2}
\;\le\;
A_\varepsilon-A_{\mathrm{sharp}}
\;\le\;
\avg{\dK\,|\nabla v_{\mathrm{sharp}}+\be|^2}. \;}
\label{eq:sandwich}
\end{equation}
Two things are worth emphasizing. First, both bounds are integrals over
$\Band$ \emph{only}, because $\dK$ vanishes elsewhere. Second, the left-hand
bound involves only $v_\varepsilon$, the corrector we already compute --- so
\eqref{eq:sandwich} is an \emph{a posteriori} error bar evaluable from the run
itself, not a statement about an unreachable limit. The derivation is four lines
and uses nothing beyond \eqref{eq:dirichlet}.

\subsection{Size of the discrepancy}

To turn \eqref{eq:sandwich} into a rate, we need the size of
$\phii^\varepsilon-\chiI$. The equilibrium profile of this model's double well
is logistic in the signed distance $s$ to $\Gam$ (positive into ice),
\begin{equation}
\phii^{\varepsilon}=\sigma(s/\varepsilon),
\qquad \sigma(u)=\frac{1}{1+e^{-u}},
\qquad \sigma(-u)=1-\sigma(u),
\label{eq:profile}
\end{equation}
and $\chiI$ is the Heaviside step $H(s)$. Integrating across the band in a
tubular neighbourhood of $\Gam$, with
$\int_{0}^{\infty}\!\big(1-\sigma(u)\big)\,du=\ln 2$,
\begin{equation}
\big\|\phii^{\varepsilon}-\chiI\big\|_{L^1(\Om)}
=\varepsilon\,|\Gam|\!\int_{-\infty}^{\infty}\!\big|\sigma(u)-H(u)\big|\,du
\;+\;\mathcal O(\varepsilon^2)
=2\ln\!2\;\varepsilon\,|\Gam|+\mathcal O(\varepsilon^2),
\label{eq:L1}
\end{equation}
the $\mathcal O(\varepsilon^2)$ coming from the curvature of $\Gam$. Feeding
\eqref{eq:dK} and \eqref{eq:L1} into \eqref{eq:sandwich},
\begin{equation}
\big|A_\varepsilon-A_{\mathrm{sharp}}\big|
\;\le\;
2\ln\!2\;(\Ki-\Ka)\,M^2\;\frac{|\Gam|}{|\Om|}\;\varepsilon
\;+\;\mathcal O(\varepsilon^2),
\qquad
M^2=\max_{\Band}|\nabla v+\be|^2 .
\label{eq:rate}
\end{equation}
This is an explicit first-order bound with a computable prefactor. Note that
$|\Gam|/|\Om|$ is the specific surface area, which the solver already reports,
so the error budget scales with a quantity we measure: at fixed $\varepsilon$,
finer microstructure means larger error, which is the correct and useful
statement of the resolution requirement.

\begin{remark}[On H-convergence]
Letting $\varepsilon\to0$ in \eqref{eq:rate} recovers
$\Keff^{\varepsilon}\to\Keff^{\mathrm{sharp}}$, which is also the standard
H-convergence result for divergence-form elliptic operators whose coefficient
converges in $L^p$. We no longer need to invoke it: for the arithmetic law the
elementary argument above gives the same conclusion \emph{plus} a constant. The
general result remains the right citation for a non-affine interpolation, where
\eqref{eq:dK} is an inequality rather than an identity.
\end{remark}

\section{The first-order correction is a surface excess}
\label{sec:excess}

Estimate \eqref{eq:rate} bounds the error but does not say what sets its sign or
how to reduce it. Both follow from looking at what the band presents to the
surrounding field. A thin graded layer of thickness $\mathcal O(\varepsilon)$ is
equivalent, to first order in $\varepsilon$, to a sharp interface carrying two
\emph{Gibbs surface excesses} per unit area --- one tangential, one normal:
\begin{equation}
\Sigma_t=\int_{-\infty}^{\infty}\!\Big[K(\phii^\varepsilon)-K_\star\Big]ds,
\qquad
\Sigma_n=\int_{-\infty}^{\infty}\!\left[\frac{1}{K(\phii^\varepsilon)}-\frac{1}{K_\star}\right]ds .
\label{eq:excess}
\end{equation}
Tangential transport crosses the layer's sublayers in \emph{parallel}, so it
sees an excess conductance $\Sigma_t$; normal transport crosses them in
\emph{series}, so it sees an excess resistance $\Sigma_n$. These two numbers are
the only $\mathcal O(\varepsilon)$ information the band transmits, so the
first-order error in $\Keff$ is a functional of $(\Sigma_t,\Sigma_n)$ alone. Both
are computable in closed form for the profile \eqref{eq:profile}.

\paragraph{The tangential excess vanishes identically.}
By affinity of $K$ and the antisymmetry $\sigma(-u)=1-\sigma(u)$, the integrand
$\sigma(u)-H(u)$ is odd, so
\begin{equation}
\Sigma_t=(\Ki-\Ka)\,\varepsilon\!\int_{-\infty}^{\infty}\!\big[\sigma(u)-H(u)\big]du=0 .
\label{eq:sigmat}
\end{equation}
This is exact, not first order: the arithmetic interpolation moves no ice across
the interface, so the phase-field cell has \emph{exactly} the sharp cell's ice
fraction, and the parallel (Voigt) response is reproduced with no error at any
$\varepsilon$. The leading correction is $\mathcal O(\varepsilon^2\kappa)$ from
interface curvature $\kappa$.

\paragraph{The normal excess does not.}
Here $1/K$ is a convex function of $\phii$, so the graded layer is a
short circuit across what should be a series resistance. Substituting
$p=\sigma(u)$, $du=dp/[p(1-p)]$ and integrating by partial fractions,
\begin{equation}
\Sigma_n
=\varepsilon\!\left[\int_{0}^{1/2}\!\!+\int_{1/2}^{1}\right]\!\!
\left[\frac{1}{\Ka+(\Ki-\Ka)p}-\frac{1}{K_\star}\right]\frac{dp}{p(1-p)}
=-\,\varepsilon\left(\frac{1}{\Ka}-\frac{1}{\Ki}\right)\ln\frac{\Ki}{\Ka}\;<0 .
\label{eq:sigman}
\end{equation}
The excess resistance is \emph{negative}: the phase-field cell is less resistive
than the sharp one, so $\Keff$ is biased \emph{high} in the direction normal to
the interface, first order in $\varepsilon$, with no free constants. For
$\Ki=2.29$, $\Ka=0.02\ \mathrm{W\,m^{-1}K^{-1}}$,
\begin{equation}
\left(\frac{1}{\Ka}-\frac{1}{\Ki}\right)\ln\frac{\Ki}{\Ka}
= 49.5633\times4.74048 = 234.959\ \ \mathrm{m\,K\,W^{-1}} .
\label{eq:const}
\end{equation}
Equations \eqref{eq:sigmat}--\eqref{eq:sigman} are the quantitative version of
the familiar warning that an arithmetic mixture rule is right in parallel and
wrong in series. They say how wrong, and in which direction.

\begin{corollary}[An interpolation with no first-order bias]
\label{cor:tensor}
The harmonic law $K_{\mathrm{harm}}(\phibar)=\big(\phii/\Ki+\phia/\Ka\big)^{-1}$
has $1/K_{\mathrm{harm}}$ affine in $\phii$, so the same antisymmetry argument
that gives \eqref{eq:sigmat} gives $\Sigma_n=0$ --- but then $\Sigma_t\neq0$.
Neither scalar law can null both. The tensorial law
\begin{equation}
\mathbf K(\phibar)
=K_{\mathrm{arith}}(\phibar)\,\big(\bI-\bn\otimes\bn\big)
+K_{\mathrm{harm}}(\phibar)\,\big(\bn\otimes\bn\big),
\qquad \bn=\frac{\nabla\phii}{|\nabla\phii|},
\label{eq:tensorK}
\end{equation}
uses each law in the direction where its excess vanishes, so
$\Sigma_t=\Sigma_n=0$ and the bias is $\mathcal O(\varepsilon^2)$. Both branches
satisfy the endpoint conditions \eqref{eq:endpoints}, so
Lemma~\ref{lem:identity} and Proposition~\ref{prop:exact} apply to
\eqref{eq:tensorK} unchanged.
\end{corollary}

\noindent
Corollary~\ref{cor:tensor} is available because the conductivity cell problem is
\emph{pure post-processing} on a frozen $\phii$ field, decoupled from the phase
evolution. The interpolation used for $\Keff$ therefore need not be the Eq.~13
used in the evolution equations, and may be chosen for accuracy.

\section{Worked example: the laminate, in closed form at every $\varepsilon$}
\label{sec:laminate}

For a layered two-phase medium every quantity above is elementary, which makes
this the natural verification case --- and the one the solver already ships as
its analytic gate.

Take a cell of period $L$ containing an ice layer of thickness $\phibar L$, with
$n_\Gam=2$ interfaces per period (under periodic boundary conditions the seam is
a real interface and must be resolved like the other). Transport along the
layers is a parallel combination and transport across them a series
combination, so $\Keff$ is diagonal with
$k_\parallel=\avg{K}$ and $k_\perp=\avg{K^{-1}}^{-1}$. Using
\eqref{eq:sigmat} and \eqref{eq:sigman},
\begin{align}
k_\parallel^{\varepsilon}
&=\Ka+(\Ki-\Ka)\,\phibar
&&\text{exact for every }\varepsilon,
\label{eq:lampar}\\[4pt]
\avg{K^{-1}}_{\varepsilon}
&=\underbrace{\frac{\phibar}{\Ki}+\frac{1-\phibar}{\Ka}}_{\text{sharp}}
\;-\;\frac{n_\Gam\varepsilon}{L}\left(\frac{1}{\Ka}-\frac{1}{\Ki}\right)\ln\frac{\Ki}{\Ka},
&&k_\perp^{\varepsilon}=\frac{1}{\avg{K^{-1}}_{\varepsilon}} .
\label{eq:lamperp}
\end{align}
Equation \eqref{eq:lamperp} is not merely a first-order truncation. The only
approximation in it is the overlap of the exponential tails of neighbouring
interfaces, which is $\mathcal O(e^{-L/4\varepsilon})$; numerically the formula
agrees with direct quadrature of $\avg{1/K(\phii^\varepsilon)}$ to $5\times
10^{-5}$ relative at $\varepsilon=L/50$ and to machine precision by
$\varepsilon=L/128$. It is, for practical purposes, exact.

\begin{center}
\begin{tabular}{@{}lrrr@{}}
\toprule
$\varepsilon$ & $\avg{K^{-1}}_{\varepsilon}$ & $k_\perp^{\varepsilon}$ & bias vs.\ $k_\perp^{\mathrm{sharp}}$ \\
\midrule
$L/50$  & 15.8200 & 0.063208 & $+59.4\%$ \\
$L/64$  & 17.8759 & 0.055941 & $+41.1\%$ \\
$L/128$ & 21.5471 & 0.046410 & $+17.0\%$ \\
$L/256$ & 23.3827 & 0.042767 & $+7.8\%$ \\
$L/512$ & 24.3005 & 0.041151 & $+3.8\%$ \\
\midrule
sharp   & 25.2183 & 0.039654 & --- \\
\bottomrule
\end{tabular}
\end{center}
{\small Predicted ladder at $\Ki/\Ka=2.29/0.02=114.5$, $\phibar=0.5$,
$n_\Gam=2$. Over the same ladder $k_\parallel=1.155000$ at every rung, to all
digits.}

\paragraph{How to plot this --- and the trap in plotting it wrong.}
Equation \eqref{eq:lamperp} is \emph{linear in $\varepsilon$ in resistivity} and
a hyperbola in conductivity. Plot and fit $1/k_\perp$ against $\varepsilon$, not
$k_\perp$. Done that way the test has two independently predicted numbers and no
fitted ones:
\begin{equation}
\text{intercept}=\frac{\phibar}{\Ki}+\frac{1-\phibar}{\Ka},
\qquad
\text{slope}=\frac{n_\Gam}{L}\left(\frac{1}{\Ka}-\frac{1}{\Ki}\right)\ln\frac{\Ki}{\Ka}
=\frac{2\times234.959}{L} .
\label{eq:fit}
\end{equation}
Plotted as $k_\perp$ versus $\varepsilon$ the same data is a curve whose
agreement can only be eyeballed.

\paragraph{Two further cautions when running it.}
Evaluate \eqref{eq:fit} at the \emph{measured} $\phibar$ reported by the solver,
not at the nominal $0.5$: the discrete field is a spline fit to nodal values and
its mean lands near but not on $0.5$, and feeding the nominal value in
manufactures an apparent solver error out of the initial condition's quadrature
error. And refine the mesh \emph{with} $\varepsilon$, holding elements per band
fixed; at fixed resolution an $\varepsilon$ ladder confounds the interface bias
measured here with ordinary mesh convergence, and the fitted slope means
nothing.

\section{Validity at finite interface width}

Collecting the results: \S4 establishes that the sharp and phase-field
formulations are the same problem, exactly, by Step~1 (reversible) and
Lemma~\ref{lem:identity} (an algebraic identity). \S\ref{sec:diff} bounds the
residual finite-$\varepsilon$ error a posteriori and at first order with a
computable constant. \S\ref{sec:excess} identifies that error as a surface
excess, shows the tangential part is exactly zero and the normal part has a
closed form, and names an interpolation \eqref{eq:tensorK} that removes both.
\S\ref{sec:laminate} gives a case where all of it is checkable in closed form.

\begin{remark}[Interpolation dependence, quantified]
The arithmetic interpolation reproduces the exact conductance for transport
\emph{tangential} to $\Gam$ --- exactly, at every $\varepsilon$, by
\eqref{eq:sigmat} --- and overestimates it for transport \emph{normal} to
$\Gam$, by \eqref{eq:sigman}. Because $\Ki/\Ka\approx114$ the asymmetry is
severe: the laminate ladder above shows a $+17\%$ bias in $k_\perp$ at
$\varepsilon=L/128$. In a real microstructure the normal and tangential
directions are mixed by the geometry, so the net bias is smaller than the
laminate's worst case but retains its sign, and it scales with the specific
surface area $|\Gam|/|\Om|$ as \eqref{eq:rate} shows.
\end{remark}

\begin{remark}[How this is defended in practice]
The claim to defend is no longer a limit theorem but the size of a known bias,
which is a much easier thing to demonstrate. The committed verification is the
$\varepsilon$ ladder of \S\ref{sec:laminate} on the periodic ice-slab cell,
checking three predictions at once: (i) $k_\parallel$ flat across the ladder to
the discretization error, confirming $\Sigma_t=0$; (ii) $1/k_\perp$ linear in
$\varepsilon$ with the intercept and slope of \eqref{eq:fit}, confirming
$\Sigma_n$; (iii) off-diagonal components zero. Passing (i) and (iii) tests the
cell solver; passing (ii) tests the theory of \S\ref{sec:excess}. For a
production microstructure, either run the same ladder and extrapolate to
$\varepsilon\to0$, or verify via \eqref{eq:rate} that the residual bias at the
production $\varepsilon$ is inside the error budget.
\end{remark}

\end{document}
